\documentclass[11pt]{article}

\usepackage[margin=1in]{geometry}
\usepackage{amsmath,amssymb}
\usepackage{booktabs}
\usepackage{microtype}
\usepackage[hidelinks]{hyperref}

% Notation used throughout.
\newcommand{\Rb}{\mathbb{R}}
\newcommand{\ip}[2]{\left\langle #1,\, #2 \right\rangle}
\newcommand{\ucon}{u_{\mathrm{contrast}}}
\newcommand{\bpre}{b_{\mathrm{pre}}}

\title{How we define and quantify causal contribution}
\author{MrVLA}
\date{}

\begin{document}
\maketitle

\begin{abstract}
\noindent
A feature's \emph{causal contribution} is its signed, exactly additive share of the
\emph{contrast-centred} action logit at the policy's readout: how much of the model's
\emph{choice between actions} that feature wrote. It is measured per decision, validated by
an additivity identity, upgraded to necessity by an exact counterfactual, and turned into a
feature-level property only through held-out, confound-controlled statistics. The definition
is \emph{functional and label-free}: nothing in it refers to what a feature means, to a firing
pattern, or to a hand-assigned category. It is a quantity in the units of the policy's own
decision.
\end{abstract}

\section{Notation}
\label{sec:notation}

\begin{center}
\begin{tabular}{@{}lll@{}}
\toprule
Symbol & Type & Meaning \\
\midrule
$h$        & $\Rb^{d}$          & layer-31 residual at the action-token position \\
$r$        & $\Rb_{>0}$         & $r=\operatorname{rms}(h)=\sqrt{\operatorname{mean}(h^{2})+\varepsilon}$ \\
$g$        & $\Rb^{d}$          & final-RMSNorm gain (entrywise) \\
$u_t$      & $\Rb^{d}$          & unembedding row for action token $t$ \\
$\ucon$    & $\Rb^{d}$          & $u_t-\tfrac{1}{256}\sum_{s}u_s$, the contrast direction \\
$z$        & $\Rb_{\ge 0}^{F}$  & \textbf{SAE code}: how strongly each feature is present \\
$w_j$      & $\Rb^{d}$          & SAE decoder row $j$, unit norm: \emph{what} feature $j$ writes \\
$\ell_2$   & $\Rb_{>0}$         & the SAE's per-sample normaliser \\
$\mu,\ \bpre$ & $\Rb,\ \Rb^{d}$ & SAE mean term and pre-bias \\
$\varphi_j$& $\Rb$              & \textbf{causal contribution} of feature $j$ to this decision \\
$C_j(g)$   & $\Rb_{\ge 0}$      & feature $j$'s causal importance on task $g$ \\
$S[j,t]$   & $\Rb$              & $\ip{w_j\odot g}{u_t}$, feature $j$'s signature over the 256 bins \\
\bottomrule
\end{tabular}
\end{center}

\noindent
Throughout, $\odot$ is the entrywise product. In the runs reported here $d$ is the OpenVLA
(Llama-2 7B) width, $F=2048$ features, and the SAE is TopK with $k=100$.

\section{Why the definition is exact rather than approximate}
\label{sec:exact}

OpenVLA emits an action as seven discrete tokens, each an argmax over 256 bins, read off the
layer-31 residual:
\begin{equation}
\operatorname{logit}(t) \;=\; \operatorname{RMSNorm}(h)\cdot u_t
\;=\; \frac{1}{r}\,\ip{h \odot g}{u_t}.
\label{eq:readout}
\end{equation}
At layer 31 the readout \emph{is} the whole remaining computation, and it is a dot product ---
additive by construction. We therefore need no gradients, no path patching and no saliency
surrogate: the decomposition below is an \emph{identity}, not an approximation.

\paragraph{Scope, stated up front.} This exactness is precisely why we work at layer 31 and
only layer 31. Earlier layers would require gradient or path-patching methods and are a
separate decision.

\section{The two symbols: \texorpdfstring{$z$}{z} and \texorpdfstring{$\varphi$}{phi}}
\label{sec:zphi}

These are the two quantities the whole argument turns on, and conflating them is exactly the
error the prior work's firing-based metric makes.

\subsection{\texorpdfstring{$z$}{z} --- presence}

$z\in\Rb_{\ge0}^{F}$ is the \textbf{SAE code} for the residual $h$: the vector of feature
activations produced by the sparse autoencoder's encoder. Because the SAE is TopK with
$k=100$ over $F=2048$ features, at most $100$ entries of $z$ are nonzero on any decision;
the rest are exactly zero. The decoder then writes the residual back as
\begin{equation}
h \;\approx\; \ell_2 \sum_{j=1}^{F} z_j\, w_j \;+\; \mu\mathbf{1} \;+\; \bpre .
\label{eq:sae}
\end{equation}
So $z_j$ answers \emph{``is feature $j$ present on this decision, and how strongly?''}, and
since $\|w_j\|=1$, the product $\ell_2 z_j$ is the amount --- in residual-stream norm --- that
feature $j$ writes into $h$.

\medskip
\noindent
$z$ is what a \textbf{firing-based} metric sees. It is a statement about the feature's
\emph{activity}, and it is silent on whether that activity affects the robot.

\subsection{\texorpdfstring{$\varphi$}{phi} --- influence}

$\varphi_j$ is the \textbf{causal contribution}: the number of logits, in the action head's own
units, that feature $j$ added to the emitted action \emph{relative to the alternatives}. It is
signed --- a feature can push the model toward the emitted bin ($\varphi_j>0$) or away from it
($\varphi_j<0$) --- and the $\varphi_j$ sum exactly to the decision (Section~\ref{sec:phi}).

\medskip
\noindent
$\varphi$ is a statement about the feature's \emph{influence}. The distinction from $z$ is the
point of the whole construction, and it is visible directly in the formula, which factorises
into three parts:
\begin{equation}
\varphi_j \;=\;
\underbrace{\frac{\ell_2}{r}}_{\text{per-decision scale}}
\;\cdot\;
\underbrace{\vphantom{\frac{\ell_2}{r}} z_j}_{\text{presence}}
\;\cdot\;
\underbrace{\vphantom{\frac{\ell_2}{r}} \ip{w_j}{g \odot \ucon}}_{\text{alignment / geometry}} .
\label{eq:phi}
\end{equation}
Only the third factor is specific to \emph{this feature acting on this action}, and it is
exactly what a firing metric cannot see. Consequently the two quantities come apart, and they
come apart in both directions:
\begin{itemize}
\item \textbf{busy-but-inert:} $z_j$ large, $\ip{w_j}{g\odot\ucon}\approx 0$, hence
      $\varphi_j\approx 0$. The feature fires constantly and writes nothing into the decision.
\item \textbf{rare-but-decisive:} $z_j$ nonzero only occasionally, but with large alignment,
      so that when it does fire it moves --- and can flip --- the emitted action.
\end{itemize}

\section{Level 1: contribution}
\label{sec:phi}

Subtracting the mean logit over the 256 action tokens from \eqref{eq:readout} gives the
\emph{contrast} logit, which is the quantity that actually decides the action (a term that
lifts every bin equally cannot change the argmax):
\begin{equation}
\operatorname{logit}(t) - \frac{1}{256}\sum_{s}\operatorname{logit}(s)
\;=\; \frac{1}{r}\,\ip{h\odot g}{\ucon}.
\label{eq:contrast}
\end{equation}
Substituting the SAE decomposition \eqref{eq:sae} into \eqref{eq:contrast} and using
bilinearity of the inner product:
\begin{equation}
\frac{1}{r}\ip{h\odot g}{\ucon}
\;=\;
\sum_{j=1}^{F}\underbrace{\frac{\ell_2}{r}\, z_j \ip{w_j}{g\odot\ucon}}_{\textstyle \varphi_j}
\;+\;
\underbrace{\frac{1}{r}\ip{\mu\mathbf{1}+\bpre}{g\odot\ucon}}_{\textstyle \text{constant}} .
\label{eq:decomp}
\end{equation}
This is the definition. Two choices in it carry the argument.

\paragraph{(i) Contrast-centring.} Attributing to $\ucon$ rather than to $u_t$ means a
direction that lifts \emph{every} action logit equally receives \emph{zero} credit. Only
movement in the direction that decides \emph{which} action wins is counted, so ``confidence''
and ``choice'' are not conflated. Without this, every high-norm feature scores.

\paragraph{(ii) The alignment term.} $\ip{w_j}{g\odot\ucon}$ is what makes the measure causal
rather than activity-based --- see Section~\ref{sec:zphi}.

\paragraph{Implementation note.} A naive $\varphi = z\ip{w}{u}$ gets two things wrong: it drops
the per-sample $\ell_2$ factor (our SAE normalises each sample), and it attributes to $u_t$
instead of $\ucon$.

\subsection{The validity gate: sufficiency}

Before attributing anything we test whether the feature terms actually recover the decision.
Sufficiency is the through-origin slope of the true action margin on the reconstructed one,
\begin{equation}
S \;=\; \frac{\sum_i c_i x_i}{\sum_i c_i^{2}},
\label{eq:suff}
\end{equation}
which decomposes into \emph{features $+$ bias $+$ error summing to $1$ by identity}.

Measured (goal suite, $k=100$, layer 31): $S = \mathbf{0.936}$, of which features alone
$=0.531$, a constant default-action bias ($\mu+\bpre$) $=0.405$, and error $=0.064$, against a
pre-registered threshold of $0.80$. The additivity canary --- do the frozen-$r$ per-feature
$\varphi_j$ sum back to the true logit --- returns correlation $1.0000$ with mean absolute
discrepancy $2.7\times10^{-14}$. The arithmetic is exact, not merely close.

The constant term is itself a small finding: roughly $40\%$ of the action margin is a fixed
lean toward a rest pose, present on every decision and belonging to no feature.

\paragraph{Recorded metric change.} This gate replaced an earlier full-residual argmax
re-decode that stalled at $0.72$ (further epochs moved it $0.70\!\to\!0.72$; raising sparsity
$k=100\!\to\!256$ reached only $0.76$ on a shallow slope). The diagnosis: the old gate punished
the SAE for reconstructing the entire dense last-layer residual --- information for the whole
$32$k vocabulary, nearly all of it action-irrelevant --- when the estimand we care about is the
\emph{action} margin. The old number is still reported alongside, and the change was recorded
before the new value was read.

\section{Level 2: necessity, by exact counterfactual}
\label{sec:necessity}

Contribution is not necessity: a large $\varphi_j$ on a decision with a large top-2 margin
changes nothing about what the robot does. At layer 31 the counterfactual can be computed
exactly. Writing the unnormalised logits $L_t = \ip{h\odot g}{u_t}$ and the signature matrix
$S[j,t] = \ip{w_j\odot g}{u_t}$, removing a set of features gives
\begin{equation}
L'_t \;=\; L_t \;-\; \sum_{k} \mathrm{coeff}_k\, S[k,t].
\label{eq:cf}
\end{equation}
Because $r>0$ rescales all $256$ bins equally, \emph{it drops out of the argmax}. Flips are
therefore exact and require \textbf{no model forward pass at all} --- every feature on every
stored decision. A top-2 margin bound prunes decisions that provably cannot flip, and the
pruned count is reported so the saving is auditable. We report:
\begin{itemize}
\item the \textbf{flip rate}: the fraction of decisions on which removing the feature changes
      the emitted bin, conditioned on the feature actually firing;
\item the \textbf{signed bin shift}: the direction and size of the change in bin \emph{index},
      which is interpretable because bins are the discretised
      $dx,dy,dz,d\mathrm{roll},d\mathrm{pitch},d\mathrm{yaw},\mathrm{gripper}$ command. The
      counterfactual is thus in units of robot motion.
\end{itemize}

\paragraph{Two ablation semantics, both reported.} These are different interventions, and the
gap between them is a real caveat rather than noise:

\begin{center}
\begin{tabular}{@{}lll@{}}
\toprule
 & $\mathrm{coeff}_j$ & what it removes \\
\midrule
\textsc{projection} & $\ip{h}{w_j}$ &
  everything along that direction: coded amount, SAE reconstruction \\
 & & error, mean term, leakage from correlated features. \emph{This is what} \\
 & & \emph{the closed-loop rollout hook does.} \\
\addlinespace
\textsc{coded} & $\ell_2 z_j$ &
  only what attribution says the feature wrote. \emph{This is what $\varphi$} \\
 & & \emph{describes.} \\
\bottomrule
\end{tabular}
\end{center}

\noindent
Their difference measures how much of a rollout ablation's effect comes from structure the SAE
never attributed to that feature --- i.e.\ how far a behavioural ablation result can honestly
be read as evidence about a \emph{coded} feature.

We also report \textbf{coalition non-orthogonality}. Decoder rows are unit-norm but not
mutually orthogonal, so the projection hook $h - (hV^{\!\top})V$ over-subtracts for correlated
coalitions. We reproduce the hook's actual formula rather than a corrected projection, because
the point is to model the experiment that was really run, and we quantify the distortion
instead of leaving it implicit.

\paragraph{Scope.} This is the exact \emph{direct} effect on one decode slot. In a rollout,
ablation also changes what slot $s{+}1$ attends to and what the next timestep observes, so the
behavioural effect is larger. Read these as a \textbf{lower bound} on behavioural impact, and
as the precise answer to ``did this feature decide this action''.

\section{Level 3: behaviour}
\label{sec:behaviour}

Closed-loop coalition ablation is the top of the ladder and the head-to-head against the
original paper's ranking: same rollout protocol, our label-free breadth ranking versus their
firing-based classifier. Success is logged \emph{per task}, on identical initial states, across
five conditions --- baseline, general coalition, specialist coalition, random matched
coalition, firing-ranked coalition --- with paired McNemar tests, damage intervals and a stated
minimum detectable effect, so that a null is separable from ``underpowered''.

The prediction that distinguishes the axis: removing a \textbf{general} coalition should damage
\emph{many} tasks (a high damage participation ratio); removing a \textbf{specialist} coalition
should damage \emph{few}.

\emph{Status: built and unit-tested, not yet run --- pending GPU allocation.}

\section{From per-decision \texorpdfstring{$\varphi$}{phi} to a property of a feature}
\label{sec:aggregation}

\begin{align}
C_j(g) &\;=\; \operatorname{mean}_{\,\text{decisions in task } g} \left| \varphi_j \right|,
\label{eq:cjg}\\[2pt]
\mathrm{PR}_j &\;=\; \frac{\bigl(\sum_g C_j(g)\bigr)^{2}}{\sum_g C_j(g)^{2}} .
\label{eq:pr}
\end{align}
The participation ratio \eqref{eq:pr} is \emph{scale-free}, so it measures the \emph{breadth} of
causal influence and not its strength: $\mathrm{PR}_j=1$ means the feature's influence lives
in a single task, $\mathrm{PR}_j=G$ means it is spread evenly across all of them. Measured on
goal: mean $\mathrm{PR}=6.05$, $p_{10}=2.65$, $p_{90}=9.91$ --- a genuine
specialist$\rightarrow$generalist spectrum rather than a two-cluster artefact.

\paragraph{A second, orthogonal breadth axis, free from the same data.} The same $|\varphi|$
mass resolved over the \emph{seven action slots}: how many of the robot's degrees of freedom a
feature actually drives. Two traps are handled explicitly. First, $\|\ucon\|$ depends on where
the emitted bin sits in the ordered range, so the gripper (near-binary, extreme bins) inflates
every feature's absolute score --- \emph{shares} are comparable across slots, absolutes are
not. Second, the gripper command is constant for most of an episode, so all statistics are
reported both on all decisions and on transitions only.

\section{What makes the number defensible}
\label{sec:defensible}

The raw statistic is \textbf{not} the claim, and we say so before anyone finds it: raw breadth
is activity-entangled, with $\operatorname{corr}(\mathrm{PR}, \text{base firing rate}) = +0.82$.
Reporting raw $\mathrm{PR}$ would be reporting firing rate with extra steps. The claim is the
confound-controlled, held-out version:

\begin{quote}
Leave one task out. Recompute \emph{both} breadth and total causal magnitude on the remaining
$G-1$ tasks --- recomputing both is what keeps the held-out task out of the confound controls
as well as out of the predictor. Then rank-partial the training breadth against the
\emph{held-out} task's causal importance, controlling \emph{both} total causal magnitude
$\sum_g C_j(g)$ and base firing rate.
\end{quote}

Measured on goal, over $446{,}096$ decisions, $10$ tasks and $2048$ features:
\begin{equation*}
\rho_{\text{partial}\,\mid\,\text{both}} \;=\; \mathbf{+0.493},
\qquad \text{positive in } 10/10 \text{ folds},
\qquad \min = +0.399, \quad \mathrm{sd} = 0.053 .
\end{equation*}

\paragraph{Supporting machinery, all implemented.}
\begin{itemize}
\item \textbf{Concentration.} Effective feature count $n_{\mathrm{eff}}$, Lorenz shares
      (top $1\%$, $10$, $50$, $100$), Gini, and cross-task top-$N$ overlap expressed as a
      \emph{ratio to chance} --- each computed against a base-firing-rate ranking as a control.
      If causal mass concentrates no more than firing rate does, ``concentration'' is an
      activity statement rather than a causal one. This attaches numbers to the reply to
      ``of course hundreds of features influence the action'': the claim is not influence, it is
      the \emph{concentration and cross-task reproducibility} of influence.
\item \textbf{Permutation floors}, with the honest note that the negative control originally
      prescribed is a \emph{no-op}: permuting task labels within a feature leaves the
      leave-one-task-out statistic unchanged, because the procedure already averages over every
      held-out task. Two floors that do test something replaced it.
\item \textbf{Attenuation correction.} For a positive result the correction is free, since
      attenuation only shrinks an observed correlation, so $r_{\mathrm{obs}}/\sqrt{r_{xx}}$ is a
      lower bound on the truth even with the other reliability unknown. Split-half breadth
      reliability is built but \emph{not yet reported}; at an illustrative $r_{xx}=0.72$ the
      floor would be $\ge +0.581$.
\end{itemize}

\paragraph{The honest limit on the claim.} $\varphi$ contains $z_j$, so $\mathrm{PR}$ is not
\emph{perfectly} activity-independent, and we do not claim that it is. The defensible statement
is that \textbf{among firing features, causal influence differs from firing frequency via the
alignment term} --- and the controls above are exactly the test of that statement.

\section{One-paragraph version for the paper}
\label{sec:paragraph}

We define a feature's causal contribution to a decision as its signed, additive share of the
contrast-centred action logit at the policy's readout. Because OpenVLA's action head is a dot
product on the final residual stream, this decomposition is exact rather than approximate: the
per-feature terms sum to the true logit to numerical precision, and the feature terms recover
$93.6\%$ of the action margin. Contribution is upgraded to necessity by an exact
counterfactual --- removing a feature's decoder direction and recomputing the argmax, which the
RMSNorm scalar leaves invariant --- yielding a per-feature flip rate and a signed shift in the
commanded action, for every feature on every decision at no forward-pass cost. Feature-level
generality is then the participation ratio of causal contribution across tasks, and is claimed
only in leave-one-task-out form, controlling both total causal magnitude and base firing rate
($\rho = +0.493$, $10/10$ folds).

\end{document}
